% area: Probability Theory

\begin{definition}[Probability Space ]
  \label{def:probspace}
  The triple $(\Omega, \mathcal{F}, P)$ is a probability space if 
  \begin{enumerate}
    \item $\Omega$ is the sample space, that is some possibly abstract set. 
    \item $\mathcal{F}$ is a $\sigma$-algebra of sets - the measurable subsets of $\Omega$. 
    \item $P$ is a probability measure. 
  \end{enumerate}

  that is, $P$ follows the Kolmogorov Axioms : 

  \begin{enumerate}
    \item For any $A \in \mathcal{F}$, there exists a number $P(A) \ge 0$, the probability of $A$. 
    \item $P(\Omega) = 1$
    \item Let $\{A_n , n \ge 1\}$ be disjoint. Then 
      $$ P\left(\bigcup_{n=1}^\infty A_n \right) = \sum_{n=1}^\infty P(A_n)
  \end{enumerate}
\end{definition}

\begin{definition}[Random Variable]
  \label{def:rv}
A random variable $X$ is a measurable function from the sample space $\Omega$ to $\R$ 

$$ X : \Omega \to \R $$ 

that is, the inverse of any Borel Set in $\R$ is $\mathcal{F}$-measurable: 
$$ X^{-1} (A) = \{\omega : X(\omega) \in A \} \in \mathcal{F} \quad \forall A \in \R $$ 
\end{definition}

\begin{definition}[Seperating Class]
  \label{def:sepclass}
 Fill in HERE % TODO
\end{definition}

\begin{definition}[Metric on the Space of Probability Measures]
  \label{def:metprobmeasure}
  Given a seperating class $\mathcal{H}$ \ref{def:sepclass}: 

  \[
    d_\mathcal{H}(F,G) = \operatorname{sup}_{h \in \mathcal{H}} | \mathbb{E}(h(F)) - \mathbb{E}(h(G)) |
  \]
\end{definition}

\begin{lemma}[Stein's Lemma]
  \label{lemma:steinslemma}
  \ref{def:metprobmeasure}
  A real valued random variable $N$ has the standard Gaussian distribution if and only if for every test function 
  $f : \mathbb{R} \to \mathbb{R}$ that is differentiable with $f' \in \mathcal{L}'(\gamma)$, the expectations 
  $\mathbb{E}(Nf(N))$ and $\mathbb{E}(f'(N)) = \mathbb{E}(Nf(N))$.
\end{lemma}

\begin{proof}
  $\implies$
  If $N \sim N(0,1)$, then 
  \[
    \mathbb{E}(f'(N)) = \mathbb{E}(Nf(N))
  \]
  implies 
  \[
    \int_\mathbb{R} f'(x) \frac{1}{\sqrt{2\pi}}dx = \int_\mathbb{R} x f(x) \frac{1}{\sqrt{2\pi}}dx 
  \]
  which if we translate into the Gaussian measure 
  \[
    \int_\mathbb{R} f'(x) \gamma(dx) = \int_\mathbb{R} xf(x) \gamma(dx)
  \]
  which is just a representation of integration by parts.

  Other direction is HW : TODO

  For some distributions, they are characterized by their moments. Necessarily this means the distribution has to be characterized by this. 
  Distribution needs moments of all moments. If you figure it out what the moments, we can take $f$ to be a monomial then we get moments of higher order, since $\mathcal{H}$ is a seperating class.
\end{proof}

\begin{remark}[Stein's Heuristic]
  \label{rmk:steinsheuristic}
  Let $F$ be a random variable such that 
  \[
    \mathbb{E}(f'(F)-Ff(f)) \approx 0
  \]
  for a large class of test functions $f$ we want to say that this is close to the standard Gaussian.

  \[
    L(F) \approx N(0,1)
  \]
where $L$ denotes the Law of a random variable (cdf).
\end{remark}
\begin{theorem}[Stein's Method]
  \label{thm:steins}
  Let $\gamma$ be the standard Gaussian Measure 
  \[
    \gamma(dx) = \frac{1}{}
  \]
\end{theorem}

\begin{definition}[Stein's Equation]
  \label{def:steinsequation}
  \ref{rmk:steinsheuristic}
  Let $N \sim N(0,1)$. Let $h : \mathbb{R} \to \mathbb{R}$ be a Borel function such that 
  \[
    \mathbb{E}(|h(N)|) < \infty 
  \]
  or in other words 
  \[
    h \in \mathcal{L}(\gamma)
  \]
  The Stein Equation associated to $h$ is the ODE 
  \[
    f'(x) -xf(x) = h(x) - \mathbb{E}(h(N))
  \]
\end{definition}

\begin{lemma}[Proposition]
  \label{lemma:ODEsolutions}
  \ref{def:steinsequation}

  All solutions of Stein's Equation are of the form 

  \[
    f(x) = Ce^{x^2/2} + e^{x^2/2}\int_{-\infty}^x [h(y)-\mathbb{E}(h(N))]e^{y^2/2} dy 
  \]
  In particular, denote by 
  \[
    f_h(x) = e^{x^2/2}\int_{-\infty}^x [h(y)-\mathbb{E}(h(N))]e^{y^2/2} dy
  \]

  In particular, $f_h$ is he only solution such that 

  \[
    \lim_{x \to \infty} e^{x^2/2} f_h(x) = 0
  \]
\end{lemma}

\begin{proof}
 $ \sqrt{2\pi}\int_{-\infty}^x h(y) \frac{e^{y^2/2}}{\sqrt{2\pi}} - \mathbb{E}(h(N))\int_{-/infty}^\infty \gamma(dx) = 0$
\end{proof}

\begin{theorem}[Stein's Lemma and Metrics via Stein's Eq]
  \label{thm:steinsmetrics}
  \ref{lemma:steinslemma}
  \ref{def:steinsequation}
  Let $\mathcal{H}$ be a seperating class and let $h \in \mathcal{H}$. 

  Let $f_h$ be the solution to the Stein Equation associated with $h$. Then 

  \[
    f'_h(x)-xf_h(x) = h(x)-\mathbb{E}(h(N))
  \]

  Let $F$ be a real valued random variable, then 
  \[
    f'_h(F)-Ff_h(F) = h(F)-\mathbb{E}(h(N))
  \]

  Taking expectations on both sides 
  \[
    \mathbb{E}(h(F))-\mathbb{E}(h(N)) = \mathbb{E}(f'_h(F)-Ff_h(F))
  \]

  Taking absolute values 

  \[
   | \mathbb{E}(h(F))-\mathbb{E}(h(N))| = |\mathbb{E}(f'_h(F)-Ff_h(F))|
  \]

  Taking a $\sup$ over $\mathcal{H}$, 

  \[
    d_{\mathcal{H}}(F,N)= \sup_{h \in \mathcal{H}}  | \mathbb{E}(h(F))-\mathbb{E}(h(N))| = \sup_{h \in \mathcal{H}} |\mathbb{E}(f'_h(F)-Ff_h(F))| 
  \]
  \ref{def:metprobmeasure}
\end{theorem}

\begin{definition}[TV Metric]
  \label{def:tvmetric}
  \ref{def:metprobmeasure}
The total variation metric is is defined as follows: 
\[
  d_\operatorname{TV}(F,G) = \sup_{B \in \mathcal{B}(\mathbb{R})} | \mathbb{P}(F \in B) - \mathbb{P}(G \in B)
\]
where $\mathcal{B}(\mathbb{R})$ are the Borel sets over the real numbers. 
\end{definition}

\begin{definition}[Kolmogorov Metric]
  \label{def:kolmetric}
  \ref{def:metprobmeasure}
  The Kolmogorow Metric metrizes the space of probability distributions given the following definition: 
  \[
    d_{\operatorname{Kol}}(F,G) = \sup | P(F \le z) - P(G\le z) |
  \]
\end{definition}

\begin{lemma}[Stein Bounds for TV Metric]
  \label{lemma:propertiesoffh}
  \ref{lemma:steinslemma}
  \ref{def:steinsequation}
  \ref{thm:steins}
  \ref{def:tvmetric}
  \ref{lemma:ODEsolutions}
   We take the seperating class that defined the total variation metric : 

   \[
   \mathcal{H}_{TV} = \{\mathbb{1}_B : B \in \mathcal{B}(\mathbb{R})\}
   \]
   The proposition is as follows: 

   Let $h : \mathbb{R} \to [0,1]$ (a bigger subset than of indicator functions), be a Borel function. Then 
Then if we denote $f_h$ the solution to the Stein equation associated with $h$ then we have 
\[
  ||f_h||< \sqrt{\pi/2} \text{ and } ||f'_h||_\infty \le 2
\]

\end{lemma}

\begin{proof}
Recall 

\[
    f_h(x) = e^{x^2/2}\int_{-\infty}^x [h(y)-\mathbb{E}(h(N))]e^{y^2/2} dy
\]

Note that 
\[
  |h(y)-\mathbb{E}(h(N))|\le 1
\]

Note that :

\[
  \int_{-\infty}^x = \int_{-\infty}^\infty - \int_x^{\infty} = - \int_x^{\infty}
\]
So, 
\[
  
    |f_h(x)| \le e^{x^2/2}\int_{-\infty}^x e^{y^2/2} dy
\]
But also 

\[
  
    |f_h(x)| \le e^{x^2/2}\int_{x}^\infty e^{y^2/2} dy
\]

If $x>0$, then 
\[
  
    |f_h(x)| \le e^{x^2/2}\int_{|x|}^\infty e^{y^2/2} dy
\]

If $x<0$, then 
\[
  
  |f_h(x)| \le e^{x^2/2}\int_{-\infty}^{|x|} e^{y^2/2} dy = e^{x^2/2}\int_{|x|}^{\infty} e^{y^2/2}  
\]

We have obtained that 

\[
  
  |f_h(x)| \le e^{x^2/2}\int_{-\infty}^{|x|} e^{y^2/2} dy = e^{x^2/2}\int_{|x|}^{\infty} e^{-y^2/2} dy \quad \forall x  
\]

The function $x \mapsto e^{x^2/2}\int_{|x|}^{\infty} e^{-y^2/2} dy$ attins its maximum at $x=0$. and 

\[
  \int_0^\infty e^{y^2/2}dy = \sqrt{\pi/2}
\]

Which proves the result by taking the $\sup$.

For the derivative, 

\[
  f_h'(x) = x e^{x^2/2}\int_{-\infty}^x [h(y)-\mathbb{E}(h(N))]e^{-y^2/2}dy 
  + 0 = h(x)-\mathbb{E}(h(N))+xe^{x^2/2}\int ... dy 
\]
Now we take bounds 
\[
  |f'_h(x)| \le 1 + |x|e^{x^2/2}\int_{-\infty}^x e^{-y^2/2}
\]

\begin{lemma}[Kol Topology > Cvg Cdf Topology]
  \label{lemma:koltopcdftop}
  \ref{def:kolmetric}
  The topology induced by the Kolmogorov metric is strictly stronger than the topology induced by convergence of CDFs of random variables.
\end{lemma}

\end{proof}
